\section{RIPPLE: Global Geometry through Local Alignment}
\label{sec:method}

RIPPLE approximates the teacher's corpus-level geometry through sparse local geometries centered at each example. A cached teacher graph defines immediate neighborhoods, diffusion composes them into broader targets, and pooled candidates realize those targets without full-corpus student scoring. Figure~\ref{fig:motivation} gives the overall view.

\subsection{From Global Geometry to Local Alignment}
\label{sec:relation-selection}

Let $\mathcal D=\{x_i\}_{i=1}^{N}$ be an unlabeled corpus, and let $t_i$ and $s_i$ be the unit-normalized teacher and student embeddings. For each anchor $i$, let $p_i^T$ be a teacher relation distribution over $\mathcal D\setminus\{x_i\}$. With teacher and student cosine similarities $c_{ij}^T$ and $c_{ij}^S$, define the bounded pair distortion
\begin{equation}
\ell_{ij}=\frac{(c_{ij}^S-c_{ij}^T)^2}{4}\in[0,1].
\end{equation}
The quantity
\begin{equation}
\mathcal E
=\frac1N\sum_i\sum_{j\neq i}p_i^T(j)\ell_{ij}
\label{eq:global-distortion}
\end{equation}
formalizes the global teacher-weighted geometry we ultimately seek to preserve; it is not an objective that we compute directly.

Conventional relational KD instead observes the examples that happen to share a mini-batch. The following proposition makes the resulting mismatch precise; Appendix~\ref{app:batch-proof} gives the proof.

\begin{proposition}[Mini-batch exposure is independent of teacher relevance]
\label{prop:batch-exposure}
Let $\mathcal O_i^{(T)}$ be the examples co-occurring with anchor $i$ over $T$ independent random reshuffles into batches of size $B$. For any teacher relation distribution $p_i^T$,
\begin{equation}
\mathbb E\!\left[1-p_i^T(\mathcal O_i^{(T)})\right]
=\left(1-\frac{B-1}{N-1}\right)^T.
\label{eq:batch-exposure}
\end{equation}
\end{proposition}

Thus important and irrelevant relations have the same inclusion probability. This suggests replacing arbitrary batch-local geometry with teacher-defined local geometry. The next proposition explains when a collection of such local objectives is sufficient; its proof is in Appendix~\ref{app:local-global-proof}.

\begin{proposition}[Local-to-global geometry bound]
\label{prop:global-geometry}
For any selected relation set $C_i$ with positive teacher mass, define
\begin{equation}
\epsilon_i
=\mathbb E_{j\sim p_i^T(\cdot\mid C_i)}[\ell_{ij}],
\qquad
\delta_i=1-p_i^T(C_i),
\end{equation}
and let $\bar\epsilon=N^{-1}\sum_i\epsilon_i$ and $\bar\delta=N^{-1}\sum_i\delta_i$. Then
\begin{equation}
\boxed{\mathcal E\leq\bar\epsilon+\bar\delta.}
\label{eq:global-geometry-bound}
\end{equation}
\end{proposition}

Global preservation therefore reduces to two requirements: accurately align the selected local geometry and cover most teacher-relevant mass. Bounds for teacher-proportional sampling and the concentrated-support comparison are deferred to Appendix~\ref{app:proportional-bound} and Appendix~\ref{app:concentrated-support}.

\subsection{Teacher-Defined Multi-Scale Local Geometry}
\label{sec:graph}

We cache the teacher embeddings and form the mutual $k$-nearest-neighbor set
\begin{equation}
\mathcal M_k^T(i)
=\operatorname{kNN}_T(i)
\cap\{j:i\in\operatorname{kNN}_T(j)\}.
\end{equation}
Mutuality suppresses one-directional hub edges. Retained edges are weighted by teacher cosine similarity and row-normalized:
\begin{equation}
P^T_{ij}
=\frac{\mathbf 1[j\in\mathcal M_k^T(i)]\exp(c^T_{ij}/\tau_i)}
{\sum_{\ell\in\mathcal M_k^T(i)}\exp(c^T_{i\ell}/\tau_i)}.
\label{eq:teacher-transition}
\end{equation}
The row $P_i^T$ defines the immediate local geometry around anchor $i$. We calibrate $\tau_i$ by a common target entropy so that neighborhood sharpness adapts to local density.

One hop preserves immediate local structure; diffusion composes overlapping neighborhoods to expose broader structure. With the lazy operator $\widetilde P^T=\tfrac12(I+P^T)$, RIPPLE defines
\begin{equation}
g^T_{i,1}=P_i^T,
\qquad
g^T_{i,r}
=\left.e_i^\top(\widetilde P^T)^r\right|_{\mathcal D\setminus\{i\}},
\quad r>1.
\label{eq:diffusion-targets}
\end{equation}
These targets provide a multi-scale view of the teacher geometry without requiring a student graph. Graph construction, entropy calibration, and sparse diffusion are detailed in Appendix~\ref{app:method-details}.

\subsection{Sparse Relational Matching}
\label{sec:candidates}

Full-corpus scoring would defeat compact training. For every $g^T_{i,r}$, RIPPLE retains high-mass entries, samples additional entries in proportion to teacher mass, and adds hard and uniform negatives. Candidate sets are pooled and deduplicated across anchors, so each candidate text is encoded once while each anchor retains its own target. Let $\Omega_{i,r}$ be the realized domain and $g^S_{i,r}$ the student cosine-softmax on that domain. The relational objective is
\begin{equation}
\mathcal L_{\mathrm{rel}}
=\frac1{|B|}\sum_{i\in B}\sum_{r\in\mathcal R}
\omega_rD_{\mathrm{KL}}\!\left(
g^T_{i,r}|_{\Omega_{i,r}}
\,\middle\|\,
g^S_{i,r}|_{\Omega_{i,r}}
\right).
\label{eq:relational-loss}
\end{equation}
The candidate construction also includes an ambient comparison over the shared pool, preventing a restricted high-mass target from hiding off-support alternatives. Exact quotas, scale weights, and temperatures are implementation details given in Appendix~\ref{app:candidates}.

Anchor-wise objectives are not independent because their local geometries overlap. RIPPLE reuses overlapping pooled candidates to supervise additional rows of the teacher graph. If $\mathcal J_B$ is the set of selected row centers, $\Omega_j^{\mathrm{row}}$ is the available teacher-neighbor support for row $j$, and $\nu_B(j)$ is its normalized selection weight, then
\begin{equation}
\mathcal L_{\mathrm{row}}
=\sum_{j\in\mathcal J_B}\nu_B(j)
D_{\mathrm{KL}}\!\left(
P_j^T|_{\Omega_j^{\mathrm{row}}}
\,\middle\|\,
\operatorname{softmax}_{\Omega_j^{\mathrm{row}}}
\left(c^S_{j\cdot}/\tau_j\right)
\right).
\label{eq:row-loss}
\end{equation}
Rows with fewer than two available teacher-neighbor columns are omitted. The complete objective is
\begin{equation}
\mathcal L_{\mathrm{RIPPLE}}
=\mathcal L_{\mathrm{rel}}
+\lambda_{\mathrm{row}}\mathcal L_{\mathrm{row}}.
\label{eq:full-objective}
\end{equation}
The cached graph is used only during distillation; inference uses the student encoder alone.
